\chapter{Dataset, Data Collection, and Labeling}

\section{Engine-Model Dataset}

The project moved from public benchmark experimentation to an internally
collected egocentric dataset. The active corpus contains 40 curated recordings:
30 for training and 10 for held-out testing. It includes 22 assembly recordings
and 18 disassembly recordings from three operators. The operator split is not
subject-disjoint, so the results measure generalization to held-out recordings
within the collected operating population rather than generalization to wholly
new operators.

\begin{table}[H]
\centering
\small
\begin{tabularx}{\textwidth}{@{}p{0.42\textwidth}X@{}}
\toprule
\textbf{Dataset characteristic} & \textbf{Value} \\
\midrule
Curated recordings & 40 \\
Train / test split & 30 / 10 recordings \\
Procedure direction & 22 assembly / 18 disassembly recordings \\
Operators & 3; not subject-disjoint between train and test \\
Source video & 1920$\times$1080 at 30 fps \\
Modeled parts & 10 named parts plus background \\
Dense label count & 386,529 labeled frame positions \\
Playable library after stress generation & 46 assets \\
\bottomrule
\end{tabularx}
\caption{Summary of the active engine-model dataset.}
\label{tab:dataset-summary}
\end{table}

\section{Recording Workflow}

The recording tool uses a phone as the egocentric camera and a local browser
page for monitoring. The host provides an HTTPS page with a WebRTC live preview;
the recording is saved locally as an MP4 and then transferred into the project
library. This made it practical to collect repeated takes while preserving the
camera perspective that an eventual wearable system would receive.

\section{Labeling Workflow}

The labeling application is a FastAPI browser tool for reviewing a recording,
scrubbing through the video, and marking the frame at which a part installation
or removal is complete. Number keys provide a fast labeling path, while the
interface records whether the procedure is assembly or disassembly. Claims,
proxies, and resumable writes allow several people to work through a library
without overwriting one another.

\begin{figure}[H]
\centering
\includegraphics[width=0.94\textwidth]{labeler_loaded.png}
\caption{Labeling application with a recording loaded. The operator can review
the video and mark procedure completion events against the part list.}
\label{fig:labeler}
\end{figure}

The tool produces a completion-label CSV, a segment CSV, and a JSON
representation. These artifacts are converted into the frame-wise labels used
for feature extraction, training, and replay visualization. Labeling the
completion event rather than only naming a clip makes the data useful for
temporal segmentation and step-transition evaluation.

\section{Label Semantics and Derived Training Data}

Each completion mark represents the point at which the annotator judged that a
part installation or removal had finished. The downstream preparation scripts
turn these sparse transition events into dense frame-wise targets: frames before
and after a transition receive the corresponding procedure-state label, and
frames in which no modeled action is active receive the background label. This
conversion is what allows the temporal model to learn both the action interval
and the boundary between adjacent steps.

The project retained the raw completion labels, the derived segments, and the
model-ready labels as separate artifacts. Keeping these levels distinct made it
possible to correct an annotation without losing the source decision, compare
different boundary interpretations, and reproduce the same feature/label
alignment during offline and online evaluation.

\section{Dataset Preparation and Split Discipline}

The production library was organized by recording rather than by randomly
sampling individual frames. Feature extraction created one fused feature
sequence per recording, and the train/test split was applied at the recording
level. This avoids placing neighboring frames from the same take in both
partitions. The current operator population is not subject-disjoint, which is
recorded as a limitation rather than hidden by a random frame split.

\section{Assembly Step Vocabulary}

\begin{table}[H]
\centering
\small
\begin{tabularx}{\textwidth}{@{}p{0.10\textwidth}p{0.43\textwidth}X@{}}
\toprule
\textbf{Key} & \textbf{Part} & \textbf{Assembly / disassembly label IDs} \\
\midrule
1 & Main Drive Planet Gear & Install 3 / remove 5 \\
2 & Pop Control Ring Gear & Install 6 / remove 8 \\
3 & Pop Control Sun Gear & Install 9 / remove 11 \\
4 & Compressor Casing & Install 12 / remove 14 \\
5 & Exhaust & Install 15 / remove 17 \\
6 & Exhaust Casing & Install 18 / remove 20 \\
7 & Cowling Bracket & Install 21 / remove 23 \\
8 & Frame Subassembly & Install 24 / remove 26 \\
9 & Propeller Cone plate & Install 27 / remove 29 \\
0 & Propeller Cone Tip & Install 30 / remove 32 \\
\bottomrule
\end{tabularx}
\caption{Procedure vocabulary used by the labeling application and model
dataset.}
\label{tab:step-vocabulary}
\end{table}

\section{Stress Variants}

Normal recordings alone do not test whether the system has learned the
procedure or has memorized a familiar order. The project therefore created
stress variants including scrambled order, reversed procedure, partial
recordings, and wrong-step cases. These assets support controlled checks for
order robustness, cold starts, and behavior when the observed action is
inconsistent with the expected sequence.

The scrambled condition is particularly important because it breaks the
assumption that the next part can be inferred from a fixed nominal order. A
reversed recording tests directionality, a partial recording tests initialization
away from the beginning of a procedure, and a wrong-step recording tests how
the application communicates activity that conflicts with the expected state.
Together these variants distinguish visual recognition from simple procedure
position tracking.
